\section{Introduction}

A long-horizon agent is supposed to keep learning after pretraining by writing things down. In every
deployed system we could inspect, what it writes down is then rewritten by a language model: after a
session the model reads its own previous memory, merges it with what just happened, and emits a
shorter replacement. The original text is not kept. This summary-of-a-summary loop is the norm, not
an edge case --- it is how context compaction works in Codex, OpenHands, Aider, Cline and Gemini
CLI, how the Anthropic API compaction endpoint is documented to behave, and how Letta/MemGPT,
Graphiti's entity summaries and LangMem maintain long-term state \citep{packer2024memgpt}. Two
timescales exist: within-session compaction, triggered at between a sixteenth and nine tenths of the
context window, and cross-session consolidation, which has \emph{no bound} on how many times a given
memory is rewritten. The second is exactly the continual-learning setting.

This loop has a structure that cognitive science has studied for two decades. A chain of learners in
which each generation treats the previous generation's output as its data is an \emph{iterated
learning} chain, and such chains are pulled toward the learners' prior
\citep{griffiths2007language,kirby2007innateness}. If that applies here, it implies something
uncomfortable, because memory exists to hold what the model does \emph{not} already expect: a user's
idiosyncratic preference, a local fact about this deployment, a lesson from one specific failure. A
process that pulls memory toward the compactor's prior is worst at keeping exactly the content
memory is for.

The intuition has been stated informally. \citet{zhang2026faulty} show that continual consolidation
can make memory \emph{worse than no memory} and that in-place rewriting is the most harmful
operator, but their synthetic world deliberately uses nonce attributes with arbitrary mappings,
which removes the prior as a variable; they attribute the loss to over-generalisation.
\citet{zahn2026facts} assert that non-default constraints ``drift toward plausible defaults'' after
three rounds of self-summarisation, but publish no per-item data and never record what a constraint
was replaced \emph{with}. The knowledge-conflict literature does grade plausibility, but only for a
single step and only for real entities, and the single-step effect is weak
\citep{clasheval2024,kochelka2026weak}. Nobody has made deviation-from-prior an experimentally
manipulated variable inside a memory chain.

We do that. Our central experiment, \textsc{E-Shift}, holds the target sentence \emph{and its value}
completely fixed and varies only a one-sentence, value-free description of the work the user does,
chosen so that the very same value is unremarkable under one description and surprising under the
other. Magnitude, wording, token length, compressibility and operational consequence are therefore
identical across arms by construction; only the compactor's expectation differs. We verify the
manipulation before running any chain, by eliciting the compactor's own prior under each frame.

We also need to know that any effect comes from the model rewriting text rather than from content
simply not fitting. We therefore build an operator lattice in which each arm removes exactly one
mechanism, and in particular a prior-blind control, \textsc{Copy-Rand}, which applies the same
capacity pressure through a uniformly random choice of which source sentences to keep. It uses the
same stimuli, the same coder and the same probes, but its selection cannot depend on any prior.

Our contributions are:

\begin{itemize}\itemsep1pt
\item \textbf{A causal test.} \textsc{E-Shift} shows \ShiftDirectionSentence{} (\S\ref{sec:c1}).
\item \textbf{A graded, controlled dose--response} across \NFamiliesConf{} minimal-pair families
whose prior-fit is measured on the compactor itself with Bradley--Terry paired comparisons, and an
operator lattice that localises the effect (\S\ref{sec:c2}).
\item \textbf{A value-level account} of the dynamics: a copy/drop/replace kernel over measured
prior-rank states, estimated from a single compaction step, and read as a quasi-stationary
distribution conditional on survival, because deletion is close to absorbing (\S\ref{sec:c6}).
\item \textbf{A memory-design consequence.} Pinning high-surprise items verbatim, using a write-time
signal a deployed system can compute for free, is compared against the anti-drift prompt industry
already ships and against token-matched random and low-surprise pinning (\S\ref{sec:c7}).
\end{itemize}

Every operating-point choice was made on a held-out development split and frozen, with the
confirmatory analysis, decision rule and effect-size thresholds fixed in advance
(\texttt{FROZEN\_PROTOCOL.md}).
